\section{从 Demo 到生产力：验收框架}

本期反复出现一个隐含评测问题：机器人到底什么时候算可用？高继扬把能力阶段从视频 demo、办公室 demo、wild demo 一直讲到生产力场景。这个划分比单看模型指标更实用，因为机器人最终要面对的不是榜单，而是真实客户和真实任务。

\subsection{四级验收}

这里把访谈里的 demo 说法整理成四级验收表。表的目的不是创造新名词，而是提醒读者：机器人能力必须从“看见一次成功”走向“持续给客户创造价值”。

\begin{longtable}{p{0.22\textwidth}p{0.34\textwidth}p{0.35\textwidth}}
\toprule
阶段 & 能证明什么 & 不能证明什么 \\
\midrule
Demo in Video & 视觉上吸引人，证明某个动作可被拍出来 & 不证明稳定性、可复现性和现场适配。 \\
Demo in Office & 在公司可控环境中可现场演示 & 不证明跨地点、跨物体、跨客户流程可用。 \\
Demo in the Wild & 在客户/展会/跨国场景中可部署演示 & 不等于长期无人值守、低维护成本。 \\
生产力场景 & 客户愿意持续使用并为价值付费 & 仍需要规模化交付、维护和数据闭环验证。 \\
\bottomrule
\end{longtable}

\begin{warningbox}{视频塑造的机器人认知偏差}
普通观众对机器人的理解往往来自短视频。视频会放大成功片段，弱化失败次数、准备时间、场景限制和维护成本。因此，判断机器人公司不能只看 demo，要看周期、出货、客户复购、维护和数据回流。
\end{warningbox}

\subsection{客户现场是最终评测}

高继扬说，客户如果评价产品不好，他会第一时间去客户现场解决，并把单个问题变成一类问题的体系化解决。这种做法对应具身智能的真实评测方式：不是离线榜单一次打分，而是在客户现场反复暴露问题、解决问题、固化流程、改进产品。

\begin{importantbox}{生产力场景的三重验证}
一个机器人能力进入生产力场景，至少要通过三重验证：客户愿意付费，系统能稳定完成任务，问题能被维护体系和数据闭环持续吸收。
\end{importantbox}

\subsection{本章小结}

机器人公司从 demo 走向生产力，需要跨过可复现、可部署、可维护、可付费四道门槛。真正的模型能力，要在这些门槛中被重新定义。

